% =====================================================================
%  Curvilinear Perspective: Geometry and Practical Implementation
%  RevTeX 4-1
% =====================================================================

\documentclass[aps,prb,reprint]{revtex4-1}

%--------------------------------------------------------------------
%  Packages
%--------------------------------------------------------------------
\usepackage{amsmath,amssymb}
\usepackage{physics}
\usepackage{graphicx}
\usepackage{hyperref}

%--------------------------------------------------------------------
\begin{document}

\title{Curvilinear Perspective: From Spherical Great Circles to Circular Arcs and Vanishing Points}
\author{Generated by ChatGPT}
\date{\today}

% ------------------------  ABSTRACT  --------------------------------
\begin{abstract}
Curvilinear perspective can be understood in two layers.  First, every point in space defines a viewing direction on a sphere centered at the eye, and every three-dimensional straight line traces a great-circle arc on that sphere.  Second, the viewing sphere must be flattened onto the image plane.  A stereographic projection maps great circles exactly to Euclidean circles or straight lines, while other wide-angle maps do not generally preserve circularity.  The practical renderer studied here uses a related constructive approach: horizontal and vertical viewing angles define two circular coordinate lines, their intersection determines each projected point, and the projected edges of a cube are drawn as circular arcs constrained by their endpoints and direction-dependent vanishing points.  We develop this geometry, determine the three pairs of cube vanishing points, and give a stable numerical procedure for constructing the finite edge arcs and their extensions.
\end{abstract}
% ---------------------------------------------------------------------

\maketitle

% ====================  MAIN SECTIONS  =================================
\section{Introduction}
Linear perspective represents the image on a plane.  It preserves straightness: a straight line in three-dimensional space is projected to a straight line in the image, and parallel lines meet at a vanishing point.  This construction is natural for a moderate field of view, but it becomes increasingly distorted as the field of view approaches or exceeds a right angle.  A curvilinear perspective instead treats the directions around the viewer as the primary geometric objects.  The scene is first projected onto a sphere centered at the eye and is then unfolded into a disk.  Peripheral directions remain at finite image coordinates, while spatial lines generally appear as curved trajectories.

There are two conceptually different reasons that circles enter this construction.  The first is exact spherical geometry: a spatial line and the eye determine a plane through the center of the viewing sphere, so the line traces a great circle on that sphere.  Under stereographic projection, every such great circle becomes an ordinary circle or a straight line in the image plane.  The second is a practical rendering choice.  The accompanying \texttt{perspective\_projection} project does not use a pure stereographic map for its active hemispherical view.  Instead, it constructs an angle-calibrated grid of circular coordinate lines and then represents every projected cube edge by the circumcircle passing through its two projected endpoints and the corresponding vanishing point.  This enforces the visual grammar of five-point curvilinear perspective even though the complete image of a three-dimensional line under the point map is not analytically guaranteed to be an exact circle.

The purpose of this note is therefore both geometric and algorithmic.  We first explain why straight lines generate great-circle arcs and review the exact circle theorem for stereographic projection.  We then derive the circular-coordinate map used in practice, show how to project the cube vertices, determine the inside and outside vanishing points associated with the three cube-axis directions, and construct the visible edge arcs and guide extensions.  Finally, we discuss numerical degeneracies and a direct test of whether the circle approximation is adequate for a chosen line segment.

\section{Viewing Sphere and Angle Definitions}
Let the eye or viewpoint be
\begin{equation}
  \vb V=(V_x,V_y,V_z),
\end{equation}
and take the forward viewing direction to be $-\hat{\vb z}$.  For a world-space point $\vb P$, define
\begin{equation}
  \vb p=\vb P-\vb V,
  \qquad
  \vb d=\frac{\vb p}{\lVert\vb p\rVert}
       =(d_x,d_y,d_z).
  \label{eq:viewing-direction}
\end{equation}
The unit vector $\vb d$ is the point on the viewing sphere associated with $\vb P$.  A visible point in front of the viewer normally has $d_z<0$.

The polar angle measured from the forward direction and the azimuth in the $x$--$y$ plane are
\begin{equation}
  \zeta=\arccos(-d_z),
  \qquad
  \psi=\operatorname{atan2}(d_y,d_x).
  \label{eq:polar-angles}
\end{equation}
Thus
\begin{equation}
  d_x=\sin\zeta\cos\psi,
  \quad
  d_y=\sin\zeta\sin\psi,
  \quad
  d_z=-\cos\zeta.
\end{equation}

For the practical construction, it is convenient to introduce separate horizontal and vertical viewing angles,
\begin{equation}
  \phi=\operatorname{atan2}(d_x,|d_z|),
  \qquad
  \theta=\operatorname{atan2}(d_y,|d_z|).
  \label{eq:theta-phi}
\end{equation}
For front-facing directions, $|d_z|=-d_z$, and therefore
\begin{equation}
  \tan\phi=\frac{d_x}{-d_z},
  \qquad
  \tan\theta=\frac{d_y}{-d_z}.
  \label{eq:tangent-angles}
\end{equation}
The central forward direction has $\theta=\phi=0$, while directions tangent to the front hemisphere satisfy $|\theta|=\pi/2$ or $|\phi|=\pi/2$ in the appropriate coordinate plane.

\section{Why a Straight Line Becomes a Great-Circle Arc}
Consider a three-dimensional line
\begin{equation}
  \vb L(t)=\vb P_0+t\vb a,
  \label{eq:spatial-line}
\end{equation}
where $\vb a$ is its direction.  The line and the viewpoint determine a plane $\Pi$ containing $\vb V$.  Every viewing ray
\begin{equation}
  \vb V+s[\vb L(t)-\vb V]
\end{equation}
lies in $\Pi$.  Consequently, the normalized directions
\begin{equation}
  \vb d(t)=
  \frac{\vb L(t)-\vb V}
       {\lVert\vb L(t)-\vb V\rVert}
  \label{eq:line-directions}
\end{equation}
also lie in $\Pi$.  Because $\Pi$ passes through the center of the viewing sphere, its intersection with the sphere is a great circle.  Hence
\begin{equation}
  \text{spatial straight line}
  \longrightarrow
  \text{great-circle arc on the viewing sphere}.
  \label{eq:line-to-great-circle}
\end{equation}
This statement is independent of the later choice of sphere-to-plane projection.

The two limiting directions of Eq.~\eqref{eq:spatial-line} are obtained from $t\to\pm\infty$:
\begin{equation}
  \lim_{t\to+\infty}\vb d(t)=\hat{\vb a},
  \qquad
  \lim_{t\to-\infty}\vb d(t)=-\hat{\vb a}.
  \label{eq:line-limits}
\end{equation}
These antipodal directions become the two vanishing points of the line family.  All spatial lines parallel to $\vb a$ share the same pair because their finite offsets disappear in the limits $t\to\pm\infty$.

\section{Exact Circle Images under Stereographic Projection}
To see when the spherical arc becomes an exact planar circle, use stereographic projection from the back pole $(0,0,1)$ of the viewing sphere.  In dimensionless coordinates,
\begin{equation}
  X=\frac{d_x}{1-d_z},
  \qquad
  Y=\frac{d_y}{1-d_z}.
  \label{eq:stereographic}
\end{equation}
The forward direction $(0,0,-1)$ maps to the origin, while the back pole is sent to infinity.  The inverse map is
\begin{align}
  d_x&=\frac{2X}{1+X^2+Y^2},\\
  d_y&=\frac{2Y}{1+X^2+Y^2},\\
  d_z&=\frac{X^2+Y^2-1}{1+X^2+Y^2}.
  \label{eq:inverse-stereographic}
\end{align}

A great circle is the intersection of the sphere with a plane through its center,
\begin{equation}
  A d_x+B d_y+C d_z=0.
  \label{eq:great-circle-plane}
\end{equation}
Substituting Eq.~\eqref{eq:inverse-stereographic} gives
\begin{equation}
  C(X^2+Y^2)+2AX+2BY-C=0.
  \label{eq:stereographic-circle}
\end{equation}
For $C\neq0$, this is a Euclidean circle.  For $C=0$, it is a straight line.  Thus stereographic projection maps every great circle to a circle or a line.

This theorem also explains the straight-line degeneracy.  A great circle passing through the stereographic projection pole maps to a line, because one point of the spherical circle has been sent to infinity.  In a curvilinear perspective image, radial and central-axis line families are therefore naturally represented by straight lines.

\section{Radial Maps and the Need for a Practical Construction}
An azimuthal equidistant, or Postel, projection centered on the forward direction is
\begin{equation}
  \vb q_{\rm P}(\vb d)
  =R\zeta\frac{(d_x,d_y)}{\sqrt{d_x^2+d_y^2}}
  =R\zeta(\cos\psi,\sin\psi),
  \label{eq:postel}
\end{equation}
with the origin defined by continuity.  It maps the front-hemisphere rim $\zeta=\pi/2$ to a circle of radius
\begin{equation}
  B=\frac{\pi R}{2}.
  \label{eq:boundary-radius}
\end{equation}
The Postel map preserves angular distance along radial image directions, but it does not generally map every great circle to a Euclidean circle.

More generally, an arbitrary radial rescaling of a stereographic image does not preserve off-center circles.  Therefore the statement
\begin{equation}
  \text{``radial rescaling preserves circularity''}
\end{equation}
is not valid without additional restrictions.  Exact circle preservation follows from the special conformal geometry of stereographic projection, not from radial symmetry alone.

The active hemispherical renderer in the project adopts a constructive alternative.  It uses angular calibration along the central horizontal and vertical axes, completes those calibrations into two families of circles, and locates a projected point by intersecting one circle from each family.  Circular projected edges are then imposed by a circumcircle construction.  This produces a visually controlled curvilinear perspective while keeping all peripheral directions finite.


\section{Circular Coordinate-Line Projection}
For a direction $\vb d$, compute $\theta$ and $\phi$ from Eq.~\eqref{eq:theta-phi}.  Define the axis calibration points
\begin{equation}
  \vb p_\theta=(0,R\theta),
  \qquad
  \vb p_\phi=(R\phi,0).
  \label{eq:axis-calibration-points}
\end{equation}
These equations make angular displacement proportional to image distance along the two central axes.

Let the four cardinal points on the $180^\circ$ boundary circle be
\begin{equation}
  X_\pm=(\pm B,0),
  \qquad
  Y_\pm=(0,\pm B),
  \label{eq:cardinal-points}
\end{equation}
where $B$ is given by Eq.~\eqref{eq:boundary-radius}.  The constant-$\theta$ coordinate circle is defined as the circle through
\begin{equation}
  X_-,\quad X_+,\quad \vb p_\theta,
\end{equation}
and the constant-$\phi$ coordinate circle is the circle through
\begin{equation}
  Y_-,\quad Y_+,\quad \vb p_\phi.
\end{equation}

Writing
\begin{equation}
  a=R\theta,
  \qquad
  b=R\phi,
\end{equation}
the two circles have equations
\begin{equation}
  x^2+y^2+\frac{B^2-a^2}{a}y-B^2=0,
  \label{eq:theta-circle}
\end{equation}
\begin{equation}
  x^2+y^2+\frac{B^2-b^2}{b}x-B^2=0.
  \label{eq:phi-circle}
\end{equation}
For $a\neq0$, the first circle has center and radius
\begin{equation}
  \vb c_\theta=
  \left(0,\frac{a^2-B^2}{2a}\right),
  \qquad
  r_\theta=\frac{a^2+B^2}{2|a|},
  \label{eq:theta-circle-data}
\end{equation}
and for $b\neq0$, the second has
\begin{equation}
  \vb c_\phi=
  \left(\frac{b^2-B^2}{2b},0\right),
  \qquad
  r_\phi=\frac{b^2+B^2}{2|b|}.
  \label{eq:phi-circle-data}
\end{equation}
The planar image $F(\vb d)$ is selected from the intersections of these two circles.

The central-axis cases are evaluated directly:
\begin{align}
  \theta=\phi=0&:\quad F(\vb d)=(0,0),\\
  \theta=0&:\quad F(\vb d)=(R\phi,0),\\
  \phi=0&:\quad F(\vb d)=(0,R\theta).
  \label{eq:axis-cases}
\end{align}
For the generic case, if two circle intersections exist, the implementation selects the point inside the boundary disk $\lVert\vb q\rVert\leq B$ that is closest to the image center.  The other intersection belongs to the unwanted continuation of the two full coordinate circles.

\section{Numerical Circle--Circle Intersection}
Let two circles have centers $\vb c_1$, $\vb c_2$ and radii $r_1$, $r_2$.  Define
\begin{equation}
  \vb \Delta=\vb c_2-\vb c_1,
  \qquad
  d=\lVert\vb \Delta\rVert.
\end{equation}
There is no real intersection if
\begin{equation}
  d>r_1+r_2,
  \qquad
  d<|r_1-r_2|,
  \qquad\text{or}\qquad
  d=0.
\end{equation}
Otherwise set
\begin{equation}
  \ell=\frac{r_1^2-r_2^2+d^2}{2d},
  \qquad
  h=\sqrt{\max(0,r_1^2-\ell^2)},
\end{equation}
\begin{equation}
  \vb m=\vb c_1+\ell\frac{\vb \Delta}{d}.
\end{equation}
The intersections are
\begin{equation}
  \vb q_\pm=
  \vb m\pm
  h\frac{(-\Delta_y,\Delta_x)}{d}.
  \label{eq:circle-intersections}
\end{equation}
The projected point is chosen by the disk and branch rules above.

In floating-point code, zero tests should be replaced by scale-aware tolerances.  Near $\theta=0$ or $\phi=0$, Eq.~\eqref{eq:axis-cases} is more stable than constructing circles whose radii diverge as $1/|\theta|$ or $1/|\phi|$.

\section{Projecting a Rotated Cube}
For a cube of side length $L$, start from the eight local vertices
\begin{equation}
  \vb u_i\in
  \left\{\left(\pm\frac L2,\pm\frac L2,\pm\frac L2\right)\right\}.
\end{equation}
If $M$ is the cube's world transformation, including its rotation and translation, the world vertices are
\begin{equation}
  \vb P_i=M\vb u_i.
  \label{eq:world-vertices}
\end{equation}
Each vertex is projected by the sequence
\begin{equation}
  \vb P_i
  \longrightarrow
  \vb d_i
  \longrightarrow
  (\theta_i,\phi_i)
  \longrightarrow
  F(\vb d_i).
  \label{eq:vertex-pipeline}
\end{equation}

The three transformed cube-axis directions can be obtained from any three edges sharing a vertex.  With the vertex ordering used in the project,
\begin{align}
  \vb a_x&=\vb P_1-\vb P_0,\\
  \vb a_y&=\vb P_3-\vb P_0,\\
  \vb a_z&=\vb P_4-\vb P_0.
  \label{eq:cube-axis-directions}
\end{align}
After normalization, every one of the twelve cube edges is assigned to one of these three direction families.  All four mutually parallel edges in one family must use the same vanishing points.

\section{Determining the Cube Vanishing Points}
A vanishing point depends only on a line direction.  For a normalized cube-axis direction $\hat{\vb a}$, consider the antipodal pair
\begin{equation}
  +\hat{\vb a},
  \qquad
  -\hat{\vb a}.
\end{equation}
Choose the member pointing more nearly forward by maximizing its dot product with
\begin{equation}
  \vb f=(0,0,-1).
\end{equation}
Call the selected direction $\vb d_{\rm in}$.  Applying the same point map used for an ordinary vertex gives the inside vanishing point
\begin{equation}
  \vb v_{\rm in}=F(\vb d_{\rm in}).
  \label{eq:inside-vp}
\end{equation}

When $d_{{\rm in},z}=0$, the direction is tangent to the front hemisphere and the vanishing point lies on the boundary circle.  Defining
\begin{equation}
  \psi_a=\operatorname{atan2}
  (d_{{\rm in},y},d_{{\rm in},x}),
\end{equation}
we set
\begin{equation}
  \vb v_{\rm in}=B(\cos\psi_a,\sin\psi_a).
  \label{eq:boundary-vp}
\end{equation}
This boundary rule avoids the singular coordinate-circle construction at a $90^\circ$ direction.

The project constructs the opposite, exterior vanishing point by inversion in the boundary circle followed by an antipodal reversal:
\begin{equation}
  \boxed{
  \vb v_{\rm out}
  =-\frac{B^2}{\lVert\vb v_{\rm in}\rVert^2}
  \vb v_{\rm in}}
  \label{eq:outside-vp}
\end{equation}
whenever $\lVert\vb v_{\rm in}\rVert$ is nonzero.  If $\vb v_{\rm in}$ lies on the boundary, Eq.~\eqref{eq:outside-vp} gives the antipodal boundary point.  If $\vb v_{\rm in}$ is at the center, the inversion is singular; the associated family is drawn as radial straight lines and no finite exterior point is required.

Equation~\eqref{eq:outside-vp} is a project-specific continuation rule.  It should not be confused with the direct image of the antipodal direction under an arbitrary sphere-to-plane map.  Its practical purpose is to provide a visually consistent second endpoint for the extension arcs outside the front viewing disk.


\section{Constructing the Circular Arc of a Cube Edge}
Let a finite cube edge have projected endpoints
\begin{equation}
  \vb q_1=(x_1,y_1),
  \qquad
  \vb q_2=(x_2,y_2),
\end{equation}
and let $\vb v=(x_3,y_3)$ be the inside vanishing point of its direction family.  The supporting circle of the edge is taken to be the circumcircle through these three points.

Define
\begin{equation}
  D=2\left[
  x_1(y_2-y_3)+x_2(y_3-y_1)+x_3(y_1-y_2)
  \right].
  \label{eq:circumcircle-determinant}
\end{equation}
If $|D|$ is below a tolerance, the three points are effectively collinear and the edge is drawn as a straight segment.  Otherwise the circle center is
\begin{align}
  c_x={}&\frac{
  (x_1^2+y_1^2)(y_2-y_3)
  +(x_2^2+y_2^2)(y_3-y_1)
  +(x_3^2+y_3^2)(y_1-y_2)}{D},\\
  c_y={}&\frac{
  (x_1^2+y_1^2)(x_3-x_2)
  +(x_2^2+y_2^2)(x_1-x_3)
  +(x_3^2+y_3^2)(x_2-x_1)}{D},
  \label{eq:circumcenter}
\end{align}
and its radius is
\begin{equation}
  r_c=\sqrt{(x_1-c_x)^2+(y_1-c_y)^2}.
  \label{eq:circumradius}
\end{equation}
This construction guarantees that the full supporting circle contains both finite endpoints and the relevant vanishing point.

The visible cube edge is only the portion between $\vb q_1$ and $\vb q_2$.  The implementation normally draws the shorter of the two arcs joining these endpoints.  The vanishing point constrains the curvature of the supporting circle but generally lies on an extension rather than on the visible finite segment.

\section{Choosing and Drawing the Arc Branch}
The angular coordinate of a planar point $\vb q$ around the circumcenter is
\begin{equation}
  \alpha(\vb q)=
  \operatorname{atan2}(q_y-c_y,q_x-c_x).
\end{equation}
Let
\begin{equation}
  \alpha_1=\alpha(\vb q_1),
  \quad
  \alpha_2=\alpha(\vb q_2),
  \quad
  \alpha_v=\alpha(\vb v).
\end{equation}
There are two arcs between $\alpha_1$ and $\alpha_2$.  For a finite cube edge, choose the local or shorter branch.  For a guide extension that must pass through a vanishing point, choose the clockwise or counterclockwise branch whose angular interval contains $\alpha_v$.

Once the branch is fixed, sample it as
\begin{equation}
  \vb q(s)=\vb c+r_c
  \bigl(\cos\alpha(s),\sin\alpha(s)\bigr),
  \qquad 0\leq s\leq1,
  \label{eq:arc-sampling}
\end{equation}
where $\alpha(s)$ interpolates continuously along the selected interval, with the appropriate $2\pi$ wrap-around.

A useful adaptive segment count is
\begin{equation}
  N_{\rm seg}=\max\left[
  N_{\min},
  \left\lceil\frac{r_c|\Delta\alpha|}{\ell_{\max}}\right\rceil
  \right],
  \label{eq:adaptive-segments}
\end{equation}
where $\ell_{\max}$ is the desired maximum arc length per rendered segment.  A fixed upper cap prevents very large circumcircles from producing excessive geometry.

\section{Drawing Guide Arcs toward Both Vanishing Points}
For each axis family, the renderer draws pale extensions from the finite edge toward $\vb v_{\rm in}$ and $\vb v_{\rm out}$.  The inside and outside extensions are handled separately:
\begin{enumerate}
  \item order the two projected vertices by their distance from $\vb v_{\rm in}$;
  \item construct an arc through the nearer vertex, the farther vertex, and $\vb v_{\rm in}$ for the inward continuation;
  \item construct the corresponding outward continuation using $\vb v_{\rm out}$ and the reversed endpoint ordering;
  \item choose the arc branch that actually passes through the requested vanishing point;
  \item place guide arcs behind the solid cube edges in rendering depth.
\end{enumerate}

Because the outside vanishing point is generated by Eq.~\eqref{eq:outside-vp}, the two guide pieces need not be portions of one globally identical Euclidean circle.  The current implementation prioritizes a consistent visual continuation on both sides of the finite edge.  If a single exact supporting curve is required, the preferable procedure is to sample the full three-dimensional line, map every sample with $F$, and render the resulting polyline or a fitted spline.

\section{Why the Circumcircle Approximation Is Reasonable}
The circumcircle construction is motivated by two exact perspective constraints.  First, a spatial line traces one great circle of viewing directions.  Second, all parallel lines share the same pair of limiting directions and hence the same vanishing points.  Requiring a projected edge to lie on a circle through its two finite endpoints and the appropriate vanishing point imposes both constraints in a simple form.

For stereographic projection, the resulting circle is exact.  For the active circular-coordinate map $F$, however, the complete mapped great circle is not analytically proven to be a Euclidean circle.  The code therefore uses the circumcircle as a geometrically constrained interpolation.  Its accuracy can be tested directly.  Sample points along the original edge,
\begin{equation}
  \vb P(s)=(1-s)\vb P_1+s\vb P_2,
  \qquad 0\leq s\leq1,
\end{equation}
map them with $F$, and calculate the circle residual
\begin{equation}
  \delta(s)=
  \left|
  \lVert F(\vb P(s))-\vb c\rVert-r_c
  \right|.
  \label{eq:circle-residual}
\end{equation}
If $\max_s\delta(s)$ is small compared with a pixel or the rendered line width, the circular approximation is visually adequate.  Otherwise the sampled mapped edge should be rendered directly.

\section{Practical Implementation Sequence}
A complete cube renderer can be organized as follows.
\begin{enumerate}
  \item Transform the eight local cube vertices into world coordinates.
  \item For every vertex, compute the unit viewing direction from the eye.
  \item Evaluate $\theta$ and $\phi$, construct the two circular coordinate lines, intersect them, and choose the valid point inside the boundary disk.
  \item Obtain the three transformed cube-axis directions from representative edges.
  \item For each axis, choose the forward orientation, map it to $\vb v_{\rm in}$, and construct $\vb v_{\rm out}$ from Eq.~\eqref{eq:outside-vp}.
  \item Associate each of the twelve edges with its $x$, $y$, or $z$ direction family.
  \item For each edge, fit the circumcircle through its two projected endpoints and the corresponding inside vanishing point.
  \item Draw the finite edge using the shorter endpoint-to-endpoint arc.
  \item Draw guide extensions toward the inside and outside vanishing points using branch-aware arcs.
  \item Replace unstable or nearly collinear constructions by straight lines.
\end{enumerate}

In the repository, the relevant routines are in \texttt{js/projections/hemispherical-projection.js}.  The cube edge list and the mapping from edge pairs to axis families are stored in \texttt{js/config.js}, while the transformed world vertices are assembled in \texttt{js/utils/three-utils.js}.

\section{Numerical and Geometric Edge Cases}
Several special cases are important in an interactive implementation.
\begin{itemize}
  \item \emph{Central vanishing point:} if $\lVert\vb v_{\rm in}\rVert$ is below a tolerance, draw the direction family as radial straight lines.
  \item \emph{Boundary vanishing point:} if $|d_z|$ is below a tolerance, use Eq.~\eqref{eq:boundary-vp} rather than the generic circle-intersection map.
  \item \emph{Nearly collinear control points:} if the determinant in Eq.~\eqref{eq:circumcircle-determinant} is small relative to the coordinate scale, draw a straight segment.
  \item \emph{Circle tangency:} clamp a slightly negative value of $r_1^2-\ell^2$ to zero before taking the square root.
  \item \emph{Wrong intersection branch:} when both coordinate-circle intersections are real, enforce the image-disk condition and select the candidate nearest the center.
  \item \emph{Angular wrap-around:} normalize angles consistently and explicitly test which clockwise or counterclockwise interval contains the control point.
  \item \emph{Rendering order:} draw guide arcs behind the solid cube arcs so that the cube silhouette remains visually dominant.
\end{itemize}

\section{Conclusion}
Curvilinear perspective is most naturally understood by separating spherical geometry from planar rendering.  A spatial straight line always produces a great circle of viewing directions, and its two orientations determine a pair of vanishing points shared by all parallel lines.  Stereographic projection maps that great circle exactly to a circle or straight line.  The practical projection used in the project follows a different but closely related construction: horizontal and vertical viewing angles define intersecting circular coordinate lines, the transformed cube axes determine three vanishing-point pairs, and each finite cube edge is represented by a circumcircle constrained by its endpoints and its direction-family vanishing point.  This yields a controllable five-point-style curvilinear perspective with clear limiting cases, while also making explicit where circularity is exact and where it is an implementation-level approximation.


% =====================================================================
\end{document}
